% 5. Concurrency control. Người viết: TV5.
% File này chỉ chứa nội dung, được \input vào report/main.tex. Nhãn có tiền tố cc-, tài liệu ở concurrency.bib.
% Số đo chỉ điền sau khi chạy (xem README.md cùng thư mục); giả thuyết ghi trước khi đo.


\noindent\textit{Người viết: Họ tên TV5}

Mục này so sánh cách Oracle và Redis xử lý nhiều yêu cầu cùng ghi vào một chỗ, qua chức năng F9 (Flash sale) của ShopFlash: một số suất giá rẻ có hạn, rất nhiều khách cùng bấm mua trong vài giây. Câu hỏi cần trả lời: khi nhiều người cùng tranh một số suất có hạn, DBMS nào giữ được đúng số lượng (không bán vượt, không mất cập nhật), với chi phí hiệu năng và độ phức tạp ra sao, và trong ShopFlash nên giao việc nào cho DBMS nào.

% -----------------------------------------------------------------------------
\subsection{Lý thuyết}\label{sec:cc-ly-thuyet}

\paragraph{Vì sao cần điều khiển đồng thời.} Một DBMS phục vụ nhiều transaction cùng lúc. Nếu các thao tác của chúng đan xen tùy ý, kết quả có thể sai dù từng transaction tự nó đúng. Điều khiển đồng thời đảm bảo chữ I (Isolation) trong ACID: kết quả của nhiều transaction chạy đan xen phải giống kết quả của một thứ tự chạy tuần tự nào đó (tính khả tuần tự, \emph{serializability}) \cite{cc-silberschatz2019}.

\paragraph{Các hiện tượng bất thường.} Bảng~\ref{tab:cc-hien-tuong} liệt kê các hiện tượng thường gặp khi không cô lập đủ, kèm ví dụ trong ShopFlash.

\begin{table}[H]
\centering
\caption{Các hiện tượng bất thường và ví dụ trong ShopFlash}\label{tab:cc-hien-tuong}
\reporttablesetup
\begin{tabularx}{\linewidth}{|P{0.28\linewidth}|L|}
\hline
\rowcolor{ReportTableHead}
\reportheader{Hiện tượng} & \reportheader{Ví dụ trong ShopFlash} \\ \hline
Mất cập nhật (lost update) & Hai khách cùng đọc \texttt{sold = 99} của sản phẩm có 100 suất, cùng ghi \texttt{sold = 100}. Cả hai được báo mua thành công nhưng bộ đếm chỉ tăng 1, thực tế bán 101 suất. \\ \hline
Đọc dữ liệu chưa commit (dirty read) & Trang flash sale hiện số suất mà một transaction khác đã trừ nhưng sau đó rollback. \\ \hline
Đọc không lặp lại (non-repeatable read) & Trong cùng một transaction, hai lần đọc số suất còn lại cho hai kết quả khác nhau. \\ \hline
Bóng ma (phantom) & Đếm số đơn của một chương trình hai lần, lần sau thấy thêm đơn vừa được chèn. \\ \hline
Lệch ghi (write skew) & Hai admin cùng kiểm tra tổng suất đã phân bổ còn dưới tồn kho rồi mỗi người thêm một sản phẩm; mỗi bên ghi một dòng khác nhau nhưng tổng cuối cùng vượt tồn kho. \\ \hline
\end{tabularx}
\end{table}

\paragraph{Mức cô lập.} Chuẩn SQL định nghĩa bốn mức cô lập theo ba hiện tượng đầu (Bảng~\ref{tab:cc-muc-co-lap}). Berenson và cộng sự chỉ ra định nghĩa này chưa đủ: cần xét thêm mất cập nhật và lệch ghi, và mức \emph{snapshot isolation} mà nhiều DBMS dùng chặn được ba hiện tượng của chuẩn nhưng vẫn để lọt lệch ghi \cite{cc-berenson1995}.

\begin{table}[H]
\centering
\caption{Mức cô lập theo chuẩn SQL (\checkmark: có thể xảy ra)}\label{tab:cc-muc-co-lap}
\reporttablesetup
\begin{tabular}{|l|c|c|c|}
\hline
\rowcolor{ReportTableHead}
\reportheader{Mức cô lập} & \reportheader{Dirty read} & \reportheader{Non-repeatable read} & \reportheader{Phantom} \\ \hline
READ UNCOMMITTED & \checkmark & \checkmark & \checkmark \\ \hline
READ COMMITTED & & \checkmark & \checkmark \\ \hline
REPEATABLE READ & & & \checkmark \\ \hline
SERIALIZABLE & & & \\ \hline
\end{tabular}
\end{table}

\paragraph{Các nhóm cơ chế.}
\begin{description}[style=nextline, leftmargin=1.25cm]
  \item[Khóa bi quan.] Khóa dữ liệu trước khi đọc để ghi, người đến sau phải chờ. Khóa hai pha (2PL) cho tính khả tuần tự nhưng có thể gây bế tắc (\emph{deadlock}): hai transaction mỗi bên giữ một khóa và chờ khóa của bên kia.
  \item[Kiểm soát lạc quan.] Không khóa khi đọc; lúc ghi mới kiểm tra dữ liệu có bị người khác đổi chưa (so sánh cột \texttt{version} hoặc giá trị cũ). Bị đổi thì thử lại. Rẻ khi ít tranh chấp, tốn kém khi tranh chấp cao vì thử lại nhiều.
  \item[MVCC.] Lưu nhiều phiên bản của dữ liệu để người đọc thấy một bản chụp nhất quán; đọc không chặn ghi và ghi không chặn đọc.
  \item[Thao tác nguyên tử.] Gom ``kiểm tra rồi ghi'' vào một lệnh duy nhất, ví dụ \texttt{UPDATE \ldots{} SET sold = sold + 1 WHERE sold + 1 <= allocated}. Không còn khoảng hở giữa lúc đọc và lúc ghi nên không còn race condition.
\end{description}

\paragraph{Khi dữ liệu nằm ở hai DBMS: Saga và TCC.} Trong ShopFlash, Redis giữ suất còn Oracle ghi đơn. Không có transaction chung cho hai hệ: commit hai pha (2PC/XA) cần cả hai DBMS tham gia, trong khi Redis không hỗ trợ XA. Giải pháp là \emph{saga} \cite{cc-garcia1987}: chia việc thành chuỗi transaction cục bộ, mỗi bước có một bước bù trừ (\emph{compensating transaction}) để hoàn lại khi bước sau thất bại. Luồng mua của ShopFlash là biến thể Try--Confirm--Cancel (TCC): bước đầu chỉ giữ chỗ (reservation có hạn), sau đó xác nhận hoặc hủy.

Saga thường được gọi là ``ACD'': có A, C, D nhưng không có I \cite{cc-richardson2018}. A ở đây là tính nguyên tử theo ngữ nghĩa (bù trừ chứ không rollback thật, trạng thái dở dang vẫn lộ ra), C chỉ đúng trong từng DBMS và đúng ``sau cùng'' trên cả hai hệ, D phụ thuộc độ bền của từng DBMS. Thiếu I nên các saga chạy cùng lúc có thể gặp chính các hiện tượng ở Bảng~\ref{tab:cc-hien-tuong}. Bảng~\ref{tab:cc-saga} cho thấy cách ShopFlash dựng lại chữ I bằng chính các kỹ thuật điều khiển đồng thời.

\begin{table}[H]
\centering
\caption{Hiện tượng do saga thiếu cô lập và biện pháp trong ShopFlash}\label{tab:cc-saga}
\reporttablesetup
\begin{tabularx}{\linewidth}{|P{0.22\linewidth}|L|P{0.29\linewidth}|}
\hline
\rowcolor{ReportTableHead}
\reportheader{Hiện tượng} & \reportheader{Tình huống trong ShopFlash} & \reportheader{Biện pháp \cite{cc-richardson2018}} \\ \hline
Đọc dữ liệu dở dang & Khách khác thấy \texttt{remaining} đã bị trừ bởi một reservation còn \texttt{PENDING}, có thể sẽ được trả lại & Khóa ngữ nghĩa (\emph{semantic lock}): trạng thái \texttt{PENDING}/\texttt{CLAIMED} \\ \hline
Mất cập nhật & Khi trả suất mà ghi đè \texttt{remaining} bằng một số tính sẵn thì làm mất lượt giữ suất của người khác & Cập nhật giao hoán (\emph{commutative update}): \texttt{HINCRBY +q}, \texttt{sold = sold + q} \\ \hline
Bù trừ hai lần & Hook rollback của lệnh đặt đơn và job dọn reservation cùng trả một suất & Kiểm tra trạng thái ngay trong Lua, chỉ chuyển sang \texttt{RELEASED} một lần \\ \hline
Quyết định trên giá trị cũ & Redis báo còn suất nhưng dữ liệu Redis vừa mất rồi được dựng lại & Kiểm tra lại (\emph{reread value}): Oracle chặn bằng \texttt{sold + q <= allocated} \\ \hline
Gửi lại yêu cầu & Mất phản hồi, khách bấm đặt đơn lần nữa & Khóa idempotency \texttt{'FS-' || reservationId} \\ \hline
\end{tabularx}
\end{table}

% -----------------------------------------------------------------------------
\subsection{Cơ chế trên Oracle}\label{sec:cc-oracle}

\begin{itemize}
  \item \textbf{MVCC dựa trên undo.} Ở READ COMMITTED (mặc định), mỗi câu truy vấn thấy dữ liệu đã commit tại thời điểm câu lệnh bắt đầu. Ở SERIALIZABLE, cả transaction thấy dữ liệu tại thời điểm transaction bắt đầu cộng với thay đổi của chính nó. Oracle không có READ UNCOMMITTED nên không bao giờ có dirty read; ngoài ra có chế độ READ ONLY \cite{cc-oracle-concepts, cc-oracle-set-transaction}.
  \item \textbf{Khóa dòng ghi ngay trong block dữ liệu}, không leo thang thành khóa bảng. Câu DML đặt khóa bảng ở mức vẫn cho nhiều phiên cùng ghi các dòng khác nhau. Đọc không chặn ghi, ghi không chặn đọc; chỉ ghi chặn ghi trên cùng một dòng \cite{cc-oracle-concepts}.
  \item \textbf{\texttt{SELECT \ldots{} FOR UPDATE}} khóa trước các dòng sẽ sửa, có thêm \texttt{NOWAIT} (gặp khóa thì báo ORA-00054 ngay), \texttt{WAIT n} và \texttt{SKIP LOCKED} \cite{cc-oracle-select}.
  \item \textbf{Tính nhất quán khi ghi.} Một câu \texttt{UPDATE} phải chờ khóa sẽ kiểm tra lại điều kiện \texttt{WHERE} trên phiên bản mới nhất của dòng sau khi lấy được khóa, và chạy lại câu lệnh nếu cần \cite{cc-oracle-concepts}. Nhờ vậy \texttt{UPDATE \ldots{} WHERE sold + 1 <= allocated} vẫn đúng ở READ COMMITTED.
  \item \textbf{SERIALIZABLE của Oracle là snapshot isolation.} Khi transaction định sửa một dòng mà transaction khác đã sửa và commit sau khi nó bắt đầu, Oracle báo ORA-08177 thay vì chờ; ứng dụng phải thử lại. Mức này chặn mất cập nhật nhưng vẫn để lọt lệch ghi \cite{cc-berenson1995}.
  \item \textbf{Deadlock} được tự phát hiện: Oracle rollback \emph{câu lệnh} (không phải cả transaction) của một phiên và báo ORA-00060. Cách phòng là mọi module khóa theo cùng một thứ tự; trong ShopFlash, \texttt{PRODUCT\_INVENTORY} luôn khóa sau cùng, nhiều dòng thì theo \texttt{product\_id} tăng dần.
  \item \textbf{Ràng buộc \texttt{CHECK} không thay được điều khiển đồng thời.} Khi mất cập nhật, giá trị ghi xuống vẫn hợp lệ (\texttt{sold = 100 <= allocated}) nên \texttt{CHECK} không phát hiện ra.
\end{itemize}

Thí nghiệm hai phiên về lost update, lock timeout và deadlock đã có ở mục Transaction (thí nghiệm B1--B3); mục này không lặp lại mà đo các chiến lược dưới tải (E1, E2) và minh họa ORA-08177 (E3).

% -----------------------------------------------------------------------------
\subsection{Cơ chế trên Redis}\label{sec:cc-redis}

\begin{itemize}
  \item \textbf{Mỗi lệnh chạy nguyên tử} vì luồng chính thực thi lệnh lần lượt từng cái; các I/O thread (từ Redis 6) chỉ lo đọc, ghi mạng. Vì vậy \texttt{INCR}, \texttt{HINCRBY} tự nó đã nguyên tử. Nhưng một quy trình nhiều lệnh (\texttt{HGET} rồi \texttt{HSET}) thì không: giữa hai lệnh vẫn có client khác chen vào.
  \item \textbf{\texttt{MULTI}/\texttt{EXEC}} xếp hàng các lệnh rồi chạy liền một mạch, không lệnh nào của client khác chen vào, nhưng \textbf{không có rollback}: lỗi lúc xếp hàng làm \texttt{EXEC} bị hủy (\texttt{EXECABORT}); lỗi lúc chạy (ví dụ sai kiểu) thì các lệnh khác trong khối vẫn được thực hiện \cite{cc-redis-transactions}.
  \item \textbf{\texttt{WATCH}} là kiểm soát lạc quan: key đang theo dõi bị đổi trước \texttt{EXEC} thì \texttt{EXEC} trả về null và client phải thử lại \cite{cc-redis-transactions}.
  \item \textbf{Lua script và Redis Functions} chạy nguyên tử: server không phục vụ lệnh nào khác trong lúc script chạy \cite{cc-redis-eval, cc-redis-functions}. Hệ quả: script phải ngắn (quá \texttt{busy-reply-threshold}, mặc định 5 giây, client khác nhận lỗi \texttt{BUSY}); và script cũng không rollback, nên script giữ suất phải kiểm tra hết điều kiện rồi mới ghi.
  \item \textbf{Không có khóa nên không có deadlock.} Khóa phân tán (\texttt{SET key value NX PX}, Redlock) là mẫu do ứng dụng tự dựng \cite{cc-redis-distlock}. Kleppmann chỉ ra tiến trình giữ khóa có thể bị dừng lâu (ví dụ do GC) tới mức khóa hết hạn mà nó không biết, nên cần thêm fencing token \cite{cc-kleppmann2016, cc-kleppmann2017}. ShopFlash không dùng khóa phân tán cho việc giữ suất.
  \item \textbf{Độ bền ảnh hưởng tới reservation.} Với AOF \texttt{everysec}, khi sập máy có thể mất khoảng 1 giây ghi đã xác nhận (xem mục Backup and recovery). Vì vậy runtime flash sale phải dựng lại được từ Oracle.
\end{itemize}

% -----------------------------------------------------------------------------
\subsection{Thiết kế thí nghiệm}\label{sec:cc-thi-nghiem}

Concurrency Lab là trang quản trị trong ShopFlash, chạy cùng một kịch bản ``N người tranh M suất'' bằng mười chiến lược (Bảng~\ref{tab:cc-chien-luoc}). Backend chạy N lượt mua trên đúng C thread (tối đa C lượt đồng thời), có cổng xuất phát chung để mọi lượt bắt đầu cùng lúc. Lab ghi vào bảng \texttt{cc\_lab\_items} và key \texttt{cclab:item} riêng, không đụng dữ liệu bán hàng.

\begin{table}[H]
\centering
\caption{Các chiến lược trong Concurrency Lab và giả thuyết (ghi trước khi đo)}\label{tab:cc-chien-luoc}
\reporttablesetup
\begin{tabularx}{\linewidth}{|P{0.27\linewidth}|L|L|}
\hline
\rowcolor{ReportTableHead}
\reportheader{Chiến lược} & \reportheader{Cách làm} & \reportheader{Giả thuyết} \\ \hline
\texttt{ORACLE\_NAIVE} & \texttt{SELECT sold}, Java kiểm tra, \texttt{UPDATE SET sold = :sold + 1} & Bán vượt, mất cập nhật; \texttt{CHECK} không báo lỗi \\ \hline
\texttt{ORACLE\_PESSIMISTIC} & \texttt{SELECT \ldots{} FOR UPDATE}, kiểm tra, \texttt{UPDATE} & Đúng; các lượt xếp hàng chờ khóa \\ \hline
\texttt{ORACLE\_NOWAIT} & \texttt{SELECT \ldots{} FOR UPDATE NOWAIT} & Không bán vượt, nhưng nhiều lượt bị từ chối (ORA-00054) dù còn hàng \\ \hline
\texttt{ORACLE\_OPTIMISTIC} & Đọc \texttt{version}, \texttt{UPDATE \ldots{} WHERE version = :v}, thử lại tối đa 5 lần & Không bán vượt; thử lại nhiều; có thể bán thiếu khi hết lượt thử \\ \hline
\texttt{ORACLE\_ATOMIC} & Một câu \texttt{UPDATE \ldots{} WHERE sold + 1 <= allocated} & Đúng, không thử lại, 1 round-trip \\ \hline
\texttt{ORACLE\_SERIALIZABLE} & Như \texttt{NAIVE} nhưng ở SERIALIZABLE & Không bán vượt; nhiều ORA-08177 \\ \hline
\texttt{REDIS\_NAIVE} & \texttt{HGET remaining}, Java kiểm tra, \texttt{HSET remaining} & Bán vượt \\ \hline
\texttt{REDIS\_WATCH} & \texttt{WATCH}, \texttt{MULTI}, \texttt{HINCRBY -1}, \texttt{EXEC}; null thì thử lại & Không bán vượt; thử lại nhiều \\ \hline
\texttt{REDIS\_DECR} & \texttt{HINCRBY -1}, âm thì \texttt{HINCRBY +1} & Đúng số bán; bộ đếm có lúc âm \\ \hline
\texttt{REDIS\_LUA} & Kiểm tra và trừ trong một Lua script & Đúng, không thử lại, 1 round-trip \\ \hline
\end{tabularx}
\end{table}

Số đo mỗi lần chạy:
\begin{itemize}
  \item số lượt được báo mua thành công;
  \item số suất bán vượt, bằng $\max(0, \text{thành công} - M)$;
  \item số cập nhật bị mất, bằng số lượt thành công trừ giá trị \texttt{sold} cuối cùng;
  \item số lần thử lại và số lỗi theo mã;
  \item throughput (lượt/giây) và độ trễ p50, p95, p99.
\end{itemize}
Mỗi cấu hình chạy 1 lần khởi động không tính rồi 5 lần đo; báo cáo trung bình $\pm$ độ lệch chuẩn mẫu.

\begin{description}[style=nextline, leftmargin=1.25cm]
  \item[E1. Bán vượt suất.] Cả 10 chiến lược; N = 1000, M = 100, C = 64.
  \item[E2. Chi phí khi tranh chấp tăng.] Các chiến lược đúng (\texttt{ORACLE\_PESSIMISTIC}, \texttt{ORACLE\_OPTIMISTIC}, \texttt{ORACLE\_ATOMIC}, \texttt{REDIS\_WATCH}, \texttt{REDIS\_LUA}); N = M = 2000 để mọi lượt đều có hàng; C $\in \{1, 8, 32, 64, 128\}$. Giả thuyết: \texttt{REDIS\_LUA} có throughput cao nhất và p95 thấp nhất; \texttt{ORACLE\_ATOMIC} hơn \texttt{ORACLE\_PESSIMISTIC} vì giữ khóa ngắn hơn; hai chiến lược lạc quan thử lại nhiều dần theo C; khi C vượt kích thước pool kết nối thì throughput không tăng nữa.
  \item[E3. SERIALIZABLE trên Oracle.] Hai phiên SQL làm lại kịch bản mất cập nhật ở mức SERIALIZABLE; giả thuyết: phiên ghi sau nhận ORA-08177.
  \item[E4. Sự cố giữa Redis và Oracle.] (a) Giữ suất rồi tắt backend, bật lại, đợi qua hạn: job dọn phải trả suất. (b) Xóa key runtime của chương trình trên Redis rồi dựng lại từ Oracle: \texttt{remaining = allocated - sold}, không bán vượt. Sau mỗi bước, view \texttt{v\_inventory\_invariant\_check} không có dòng nào.
  \item[E5. Luồng giữ suất qua HTTP] (nếu còn thời gian). 500 khách cùng giữ 100 suất qua API bằng k6.
\end{description}

\subsection{Kết quả}\label{sec:cc-ket-qua}

\paragraph{Môi trường đo.}
\chuaviet{Máy, hệ điều hành, Docker; phiên bản Oracle, Redis (AOF); JDK, Spring Boot; kích thước pool kết nối Oracle và Redis.}

\paragraph{E1.}
\chuaviet{Bảng kết quả 10 chiến lược (bán vượt, mất cập nhật, thử lại, lỗi theo mã) và ảnh trang Concurrency Lab.}

\paragraph{E2.}
\chuaviet{Biểu đồ throughput và p95 theo C; bảng số liệu.}

\paragraph{E3, E4.}
\chuaviet{Các bước kèm ảnh chụp; trạng thái runtime và view bất biến sau mỗi bước.}

% -----------------------------------------------------------------------------
\subsection{So sánh}\label{sec:cc-so-sanh}

\begin{table}[H]
\centering
\caption{So sánh điều khiển đồng thời trên Oracle và Redis}\label{tab:cc-so-sanh}
\reporttablesetup
\begin{tabularx}{\linewidth}{|P{0.24\linewidth}|L|L|}
\hline
\rowcolor{ReportTableHead}
\reportheader{Tiêu chí} & \reportheader{Oracle} & \reportheader{Redis} \\ \hline
Đơn vị tranh chấp & Dòng (khóa dòng), không leo thang khóa & Cả server: mỗi lệnh hoặc script chạy một mình \\ \hline
Đọc và ghi & MVCC, đọc không chặn ghi & Không cần, mọi lệnh đã tuần tự \\ \hline
Mức cô lập & READ COMMITTED (mặc định), SERIALIZABLE (snapshot isolation), READ ONLY & Không có khái niệm mức cô lập; \texttt{MULTI}/\texttt{EXEC} chỉ đảm bảo không bị chen ngang \\ \hline
Rollback & Có, theo câu lệnh và theo transaction & Không (cả \texttt{MULTI}/\texttt{EXEC} lẫn Lua) \\ \hline
Khóa bi quan & \texttt{SELECT \ldots{} FOR UPDATE [NOWAIT | WAIT n | SKIP LOCKED]} & Không có sẵn; tự dựng khóa phân tán \\ \hline
Khóa lạc quan & Cột \texttt{version}, \texttt{UPDATE \ldots{} WHERE version = :v} & \texttt{WATCH} \\ \hline
Kiểm tra rồi ghi nguyên tử & Một câu \texttt{UPDATE \ldots{} WHERE} có điều kiện & Lua script hoặc Function \\ \hline
Deadlock & Có thể xảy ra; tự phát hiện (ORA-00060) & Không có \\ \hline
Độ bền sau khi ghi & Redo log & Tùy cấu hình AOF, RDB \\ \hline
Throughput khi tranh một dòng hoặc key (E2) & \chuaviet{} & \chuaviet{} \\ \hline
Bán vượt khi viết theo cách ngây thơ (E1) & \chuaviet{} & \chuaviet{} \\ \hline
\end{tabularx}
\end{table}

% -----------------------------------------------------------------------------
\subsection{Kết luận trong bối cảnh ShopFlash}\label{sec:cc-ket-luan}

\chuaviet{Viết sau khi có số đo; chỉ giữ những ý được số đo ủng hộ. Hướng dự kiến: giữ suất flash sale dùng Redis với Lua (E1, E2); ghi đơn, trừ kho dùng Oracle với UPDATE có điều kiện (E1); ranh giới giữa hai hệ là reservation có hạn kèm bù trừ, Oracle là chốt chặn cuối (E4); giới hạn của phép đo: một máy, Docker, một backend.}
